\section{Chirality and Helicity in the Massless Limit}

The Dirac equation is
\begin{align}
(\gamma^\mu p_\mu - m)u(p) = 0 .
\end{align}

In the chiral basis,
\begin{align}
\begin{pmatrix}
-m I_{2\times 2} & \sigma \cdot p \\
\bar{\sigma} \cdot p & -m I_{2\times 2}
\end{pmatrix}
\begin{pmatrix}
u_L \\
u_R
\end{pmatrix}
= 0 .
\end{align}

Equivalently,
\begin{align}
(\sigma \cdot p)u_R - m u_L &= 0, \\
(\bar{\sigma} \cdot p)u_L - m u_R &= 0 .
\end{align}

In the massless limit,
\begin{align}
m \to 0
\quad\Longrightarrow\quad
\begin{cases}
(p^0 - \boldsymbol{\sigma}\cdot\mathbf{p})u_R = 0, \\
(p^0 + \boldsymbol{\sigma}\cdot\mathbf{p})u_L = 0 .
\end{cases}
\end{align}

For a massless particle, \(p^0 = |\mathbf{p}|\), so
\begin{align}
\frac{\boldsymbol{\sigma}\cdot\mathbf{p}}{|\mathbf{p}|}u_R &= u_R, \\
\frac{\boldsymbol{\sigma}\cdot\mathbf{p}}{|\mathbf{p}|}u_L &= -u_L .
\end{align}

The helicity operator is
\begin{align}
\hat{h}
=
\frac{\mathbf{S}\cdot\mathbf{p}}{|\mathbf{p}|}
=
\frac{1}{2}
\frac{\boldsymbol{\sigma}\cdot\mathbf{p}}{|\mathbf{p}|}.
\end{align}

Therefore,
\begin{align}
\hat{h}u_R &= +\frac{1}{2}u_R, \\
\hat{h}u_L &= -\frac{1}{2}u_L .
\end{align}

Thus, for massless fermions, chirality states coincide with helicity states.

Conceptually one can think of this phenomenon in terms of Lorentz boosts. 
A massless particle always travels at the speed of light, 
so no Lorentz boost can reverse its direction of motion. 
Therefore, the helicity of a massless particle is Lorentz invariant and coincides with its chirality. 
In contrast, for a massive particle, a suitable Lorentz boost can reverse its direction of motion, 
changing its helicity while leaving its chirality unchanged.
